% GENERATED FILE. Edit canonical lessons, then run npm run generate:books.
% canonical-source-hash: fnv1a64:0072b0d4941684b2
% canonical-lessons: LA-C195-libum, LA-C195-acetum, LA-C195-garum, LA-C195-perna, LA-C195-dies-natalis

\chapter{Cake, Vinegar, Fish sauce, Ham, Birthday}
\label{ch:la-cake-vinegar-fish-sauce-ham-birthday}
\hlchaptermodality{195}

\begin{chapteropening}
\textbf{By the end of this chapter:} \emph{I can say a cake, vinegar, fish sauce, a ham and a birthday.}

Complete the last lesson of chapter 195: I can say a cake, vinegar, fish sauce, a ham and a birthday.
\end{chapteropening}

\section[lībum, lībī]{lībum, lībī — a cake}

\label{lesson:LA-C195-libum}



\begin{quote}
Before the new one: say the Latin for a rumour, then the Latin for history.
\end{quote}

\subsection*{You'll want to know: lībum, lībī}
\textbf{lībum, lībī} — \textquotedblleft{}a cake\textquotedblright{}.

\begin{cousinweb}
\textbf{lībum} has no sure origin; Varro tied it to \emph{lībō}, \textquotedblleft{}to pour out an offering\textquotedblright{} (Lewis \& Short), since such cakes were offered to the gods.
\end{cousinweb}

\begin{grammarlens}[title={What you've built}]
One more word for this chapter.
\end{grammarlens}

\subsection*{Your turn}
\begin{itemize}
\raggedright
  \item \emph{Say it:} \emph{lībum, lībī}
  \item \emph{Say it:} \emph{lībum, lībī}, once more
  \item \emph{Say it:} say \emph{lībum, lībī}
  \item \emph{Recall:} say the Latin for a rumour, then the Latin for history, then say \emph{lībum, lībī} again
\end{itemize}

\begin{tcolorbox}[breakable,colback=teal!4,colframe=teal!35!black,title={Before you move on}]
What does lībum, lībī mean? (\textquotedblleft{}A cake\textquotedblright{}.) Say it once more.
\end{tcolorbox}
\section[acētum, acētī]{acētum, acētī — vinegar}

\label{lesson:LA-C195-acetum}



\begin{quote}
Before the new one: say the Latin for history, then the Latin for a cake.
\end{quote}

\subsection*{You'll want to know: acētum, acētī}
\textbf{acētum, acētī} — \textquotedblleft{}vinegar\textquotedblright{}.

\begin{cousinweb}
\textbf{acētum} is the participle of \emph{aceō}, \textquotedblleft{}to be sour\textquotedblright{} (Lewis \& Short), with \emph{vīnum}, \textquotedblleft{}wine\textquotedblright{}, understood: soured wine. \emph{aceō} belongs with \emph{ācer}, \textquotedblleft{}sharp\textquotedblright{}.
\end{cousinweb}

\begin{grammarlens}[title={What you've built}]
One more word for this chapter.
\end{grammarlens}

\subsection*{Your turn}
\begin{itemize}
\raggedright
  \item \emph{Say it:} \emph{acētum, acētī}
  \item \emph{Say it:} \emph{acētum, acētī}, once more
  \item \emph{Say it:} say \emph{acētum, acētī}
  \item \emph{Recall:} say the Latin for history, then the Latin for a cake, then say \emph{acētum, acētī} again
\end{itemize}

\begin{tcolorbox}[breakable,colback=teal!4,colframe=teal!35!black,title={Before you move on}]
What does acētum, acētī mean? (\textquotedblleft{}Vinegar\textquotedblright{}.) Say it once more.
\end{tcolorbox}
\section[garum, garī]{garum, garī — fish sauce}

\label{lesson:LA-C195-garum}



\begin{quote}
Before the new one: say the Latin for a cake, then the Latin for vinegar.
\end{quote}

\subsection*{You'll want to know: garum, garī}
\textbf{garum, garī} — \textquotedblleft{}fish sauce\textquotedblright{}.

\begin{cousinweb}
\textbf{garum} was borrowed from Greek \emph{garon}, \textquotedblleft{}fish sauce\textquotedblright{} (Lewis \& Short).
\end{cousinweb}

\begin{grammarlens}[title={What you've built}]
One more word for this chapter.
\end{grammarlens}

\subsection*{Your turn}
\begin{itemize}
\raggedright
  \item \emph{Say it:} \emph{garum, garī}
  \item \emph{Say it:} \emph{garum, garī}, once more
  \item \emph{Say it:} say \emph{garum, garī}
  \item \emph{Recall:} say the Latin for a cake, then the Latin for vinegar, then say \emph{garum, garī} again
\end{itemize}

\begin{tcolorbox}[breakable,colback=teal!4,colframe=teal!35!black,title={Before you move on}]
What does garum, garī mean? (\textquotedblleft{}Fish sauce\textquotedblright{}.) Say it once more.
\end{tcolorbox}
\section[perna, pernae]{perna, pernae — a ham}

\label{lesson:LA-C195-perna}



\begin{quote}
Before the new one: say the Latin for vinegar, then the Latin for fish sauce.
\end{quote}

\subsection*{You'll want to know: perna, pernae}
\textbf{perna, pernae} — \textquotedblleft{}a ham\textquotedblright{}.

\begin{cousinweb}
\textbf{perna} is inherited, usually traced to Proto-Indo-European \emph{*persneh2}, \textquotedblleft{}heel\textquotedblright{} (EtymDB 2.1): first the leg of meat with the heel, then the ham. Gothic \emph{fairzna}, \textquotedblleft{}heel\textquotedblright{}, is a cousin.
\end{cousinweb}

\begin{grammarlens}[title={What you've built}]
One more word for this chapter.
\end{grammarlens}

\subsection*{Your turn}
\begin{itemize}
\raggedright
  \item \emph{Say it:} \emph{perna, pernae}
  \item \emph{Say it:} \emph{perna, pernae}, once more
  \item \emph{Say it:} say \emph{perna, pernae}
  \item \emph{Recall:} say the Latin for vinegar, then the Latin for fish sauce, then say \emph{perna, pernae} again
\end{itemize}

\begin{tcolorbox}[breakable,colback=teal!4,colframe=teal!35!black,title={Before you move on}]
What does perna, pernae mean? (\textquotedblleft{}A ham\textquotedblright{}.) Say it once more.
\end{tcolorbox}
\section[diēs nātālis, diēī nātālis]{diēs nātālis, diēī nātālis — a birthday}

\label{lesson:LA-C195-dies-natalis}



\begin{quote}
Before the new one: say the Latin for fish sauce, then the Latin for a ham.
\end{quote}

\subsection*{You'll want to know: diēs nātālis, diēī nātālis}
\textbf{diēs nātālis, diēī nātālis} — \textquotedblleft{}a birthday\textquotedblright{}.

\begin{cousinweb}
\textbf{diēs nātālis} is \emph{diēs}, \textquotedblleft{}day\textquotedblright{}, with \emph{nātālis}, \textquotedblleft{}of one's birth\textquotedblright{}, built on \emph{nātus}, \textquotedblleft{}born\textquotedblright{}, from \emph{nāscor}, \textquotedblleft{}to be born\textquotedblright{}. English \emph{natal} comes from the second word.
\end{cousinweb}

\begin{grammarlens}[title={What you've built}]
One more word for this chapter.
\end{grammarlens}

\subsection*{Your turn}
\begin{itemize}
\raggedright
  \item \emph{Say it:} \emph{diēs nātālis, diēī nātālis}
  \item \emph{Say it:} \emph{diēs nātālis, diēī nātālis}, once more
  \item \emph{Say it:} say \emph{diēs nātālis, diēī nātālis}
  \item \emph{Recall:} say the Latin for fish sauce, then the Latin for a ham, then say \emph{diēs nātālis, diēī nātālis} again
\end{itemize}

\begin{tcolorbox}[breakable,colback=teal!4,colframe=teal!35!black,title={Before you move on}]
What does diēs nātālis, diēī nātālis mean? (\textquotedblleft{}A birthday\textquotedblright{}.) Say it once more.
\end{tcolorbox}
